\section*{Capítulo \ref{ch:g11:vision-and-brain} --- A visão e o encéfalo}

\begin{solution}{exo:g11:vision-and-brain:1}
A luz do campo esquerdo cai sobre a metade direita de cada retina; as
fibras da metade direita do olho esquerdo cruzam no quiasma, as do olho
direito ficam à direita; juntas elas formam o trato ótico direito,
fazem escala no tálamo e chegam ao \omterm{prop:g11:vision-and-brain:pathway}{córtex visual primário} direito.
\end{solution}

\begin{solution}{exo:g11:vision-and-brain:2}
Os dois nervos óticos se encontram; as fibras da metade nasal de cada
retina cruzam para o outro lado, as da metade temporal não. Depois
dele, cada trato carrega uma metade do campo visual.
\end{solution}

\begin{solution}{exo:g11:vision-and-brain:3}
Área da cor (mundo visto em cinza), área do movimento (nenhuma
percepção de movimento), área dos rostos (rostos não reconhecidos).
\end{solution}

\begin{solution}{exo:g11:vision-and-brain:4}
Plasticidade: a capacidade de as conexões dos neurônios serem
reforçadas, enfraquecidas, formadas ou podadas conforme o uso. Período
crítico: a fase inicial da vida durante a qual a experiência molda o
córtex de modo decisivo e irreversível.
\end{solution}

\begin{solution}{exo:g11:vision-and-brain:5}
Ele se liga aos receptores de serotonina de neurônios das áreas
visuais, perturbando o fluxo de sinais e produzindo alucinações sem
nenhuma luz correspondente.
\end{solution}

\begin{solution}{exo:g11:vision-and-brain:6}
Depois do quiasma, à esquerda: o trato ótico esquerdo, a escala
talâmica esquerda ou o \omterm{prop:g11:vision-and-brain:pathway}{córtex visual primário} esquerdo.
\end{solution}

\begin{solution}{exo:g11:vision-and-brain:7}
As fibras que cruzam vêm das metades nasais das duas retinas, que veem
as metades externas (temporais) do campo: o paciente perde a metade
esquerda do campo do olho esquerdo e a metade direita do campo do olho
direito --- visão em túnel dos dois lados.
\end{solution}

\begin{solution}{exo:g11:vision-and-brain:8}
Normal: $12 + 20 + 36 + 20 = 88\%$. Privado: $2 + 3 + 5 + 30 = 40\%$,
dos quais apenas 10\% principalmente.
\end{solution}

\begin{solution}{exo:g11:vision-and-brain:9}
Tapar força o córtex a usar a entrada do olho fraco, de modo que as
conexões dele são conservadas e reforçadas durante o período crítico em
vez de podadas. Num adulto o período acabou: as conexões não se
reorganizam mais nessa escala, de modo que o olho fraco não consegue
retomar o território dele.
\end{solution}

\begin{solution}{exo:g11:vision-and-brain:10}
Os \omterm{def:g11:the-eye:photoreceptors}{cones} e as \omterm{def:g11:the-eye:photoreceptors}{opsinas} dela são normais; o que se perdeu é a etapa
cortical que interpreta as razões entre eles. Um registro elétrico da
resposta da retina a luzes coloridas seria normal nela e anormal numa
pessoa daltônica; e o déficit dela é total e adquirido, o delas parcial
e de nascença.
\end{solution}

\begin{solution}{exo:g11:vision-and-brain:11}
Anos de prática reforçaram e multiplicaram as conexões que servem os
dedos dela, aumentando o território cortical delas. Sem prática, a área
encolheria de novo ao longo de anos, como mostram os estudos com
malabarismo.
\end{solution}

\begin{solution}{exo:g11:vision-and-brain:12}
No começo ela vê luz, cores e manchas em movimento, mas não consegue
reconhecer formas, rostos nem profundidade: a retina funciona, mas o
córtex nunca construiu, durante o período crítico, as conexões que
interpretam os sinais dela. Ela vai aprender a usar a visão em parte,
devagar, mas nunca com a fluência de quem enxerga desde o nascimento.
\end{solution}

\begin{solution}{exo:g11:vision-and-brain:13}
Na fóvea cada cone tem a \omterm{def:g10:cells-common-unit:cell}{célula} ganglionar própria dele, de modo que o
centro do campo envia muito mais fibras por grau que a periferia, onde
cem \omterm{def:g11:the-eye:photoreceptors}{bastonetes} compartilham uma; o córtex distribui a área em proporção
às fibras, e portanto ao detalhe, e não aos graus de campo.
\end{solution}

\begin{solution}{exo:g11:vision-and-brain:14}
A atividade da área é a percepção da cor, venha o sinal do olho ou da
memória; a percepção é o que as áreas fazem, e não o que o olho
entrega. Uma alucinação é essa atividade disparada por uma droga, em
vez de pela intenção ou pela luz.
\end{solution}

\begin{solution}{exo:g11:vision-and-brain:15}
O olho converte a luz em sinais; quem vê é o encéfalo. Um paciente com
olhos intactos pode perder a cor, o movimento ou os rostos por uma
lesão cortical; um gatinho com olhos intactos fica funcionalmente cego
de um deles se o córtex não puder conectá-lo; e uma droga que não toca
nem no olho nem na estrutura do córtex faz uma pessoa ver o que não
está lá. O olho fornece; o encéfalo vê.
\end{solution}

\begin{solution}{pb:g11:vision-and-brain:1}
\textbf{1.} O olho esquerdo ou o nervo ótico dele, antes do quiasma:
só as fibras de um olho estão afetadas.

\textbf{2.} O trato ótico direito, o tálamo ou o \omterm{prop:g11:vision-and-brain:pathway}{córtex visual}
direito: um déficit compartilhado pelos dois olhos para o mesmo meio
campo só pode estar onde as fibras dos dois olhos para aquela metade
foram reunidas, ou seja depois do quiasma.

\textbf{3.} O próprio quiasma: as fibras que cruzam, vindas das metades
nasais das duas retinas, que veem as metades temporais do campo, foram
cortadas.

\textbf{4.} A área do movimento, depois do córtex primário. Campos e
acuidade intactos mostram que o córtex primário e tudo o que vem antes
dele estão funcionando.

\textbf{5.} B, C e D têm retinas intactas; a de A pode estar ou não.
Verifica-se examinando a retina com um oftalmoscópio e registrando a
resposta elétrica dela a clarões.

\textbf{6.} O córtex direito recebe a metade esquerda do campo.

\textbf{7.} $2^\circ$ centrais: \qty{1}{cm}; uma faixa de $2^\circ$ a
$40^\circ$: \qty{0.1}{cm}. Razão 10.

\textbf{8.} A fóvea empacota os \omterm{def:g11:the-eye:photoreceptors}{cones} densamente, cada um com a \omterm{def:g10:cells-common-unit:cell}{célula}
ganglionar própria dele, de modo que os graus centrais enviam muito
mais fibras que a periferia; o córtex lhes dá proporcionalmente mais
área.

\textbf{9.} Um pequeno ponto cego no centro do campo esquerdo, o que é
devastador para a leitura; mais para a frente, uma região cega maior na
periferia do campo esquerdo, menos notada.

\textbf{10.} Movendo os olhos, o paciente traz cada palavra para uma
parte da retina cujo córtex está intacto, usando a borda da lesão;
lento, mas possível.

\textbf{11.} $20 + 36 + 20 = 76\%$.

\textbf{12.} $2 + 3 + 5 + 30 = 40\%$, e apenas 10\% com o olho esquerdo
como entrada principal ou equivalente.

\textbf{13.} A entrada dele chega ao córtex, mas as conexões dele foram
enfraquecidas e podadas por falta de uso; as conexões do olho direito
se expandiram para o espaço.

\textbf{14.} Que a reorganização fica confinada a um período crítico do
começo da vida; o equivalente humano é a ambliopia, que só se
desenvolve na infância e é intratável depois dela.

\textbf{15.} Sem nenhum olho favorecido, não há tomada: o histograma
fica aproximadamente simétrico, mas com menos neurônios responsivos ao
todo, e a visão do gatinho fica ruim nos dois olhos --- a privação sem
competição enfraquece os dois.

\textbf{16.} O encéfalo suprime a imagem do olho desviado para evitar a
visão dupla; sem uso, as conexões dele são podadas durante o período
crítico e o olho fica permanentemente fraco. Os óculos corrigem a
ótica, mas o olho não é usado, de modo que a poda prossegue.

\textbf{17.} Cobrir o olho bom força o córtex a usar os sinais do olho
fraco, conservando e reforçando as conexões dele.

\textbf{18.} O período crítico termina por volta dos sete anos; depois
dele o córtex não redistribui mais território nessa escala, de modo que
as conexões perdidas do olho fraco não podem ser reconstruídas.

\textbf{19.} No adulto as conexões já estão estabelecidas e estáveis; a
tomada pelo outro olho é um fenômeno do córtex em desenvolvimento,
durante o período crítico, e não do córtex maduro.

\textbf{20.} A: olho ou nervo ótico; C: quiasma; B: trato direito,
tálamo ou córtex primário; D: a área do movimento, depois deles. A
plasticidade do córtex, confinada a um período crítico, faz o
tapa-olho da criança funcionar aos quatro anos e falhar aos quinze.
\end{solution}
